\subsection*{IV 20.2: Etude de la section hyperplane générique---propriétés locales (cote 25, sites 8--15)}
The section studies the generic hyperplane through the incidence construction of no.\ 1. Propositions 2.1--2.6 concern irreducibility, empty inverse images, the induced correspondence of irreducible components, and dimension; proposition 2.7 treats \emph{cycles premiers associés}. The proof refers to paragraph 3 for associated cycles under the stated flat morphism, then explicitly shortens its own route through the universal incidence family (site 10). Proposition 2.8 and its corollaries use paragraph 6 for the local properties \emph{coprof}, $(S_k)$, Cohen--Macaulay, reducedness and integrality. Corollary 2.11 adds regularity, $(R_k)$ and normality. The discussion then turns to smoothness and the corresponding geometric properties in 2.12--2.15, using the \emph{critère jacobien de lissité}, followed by \emph{multiplicité radicielle} and separability in 2.16--2.17. The handwriting changes hypotheses and the organization: a cancelled transition on site 11 reappears on site 12; the inserted hypothesis preceding 2.12 and the margin beside 2.14 contain a ramifié qualification whose initial letters remain doubtful; site 14 proposes presenting the practical corollary before its supporting result; Remarques 2.18 revisit the immersion assumption and characteristic zero. The checked pages retain several mathematical slips or gaps: the index is absent from $k(\ )$ on site 13, the argument there prints $p\leq n$ before its separate $p=n$ discussion, and both occurrences in the non-embeddedness assertion of 2.16 on site 14 read $Z_i$ without an added generic-point index. The guide records these forms without repairing their hypotheses or formulas. Sources: \href{https://grothendieck.umontpellier.fr/web/viewer.html?file=../25.pdf\#page=9}{cote 25, sites 8--15}; internal references to published EGA IV, paragraphs 3 and 6, are those explicitly written in the prenotes.
